\documentclass[10pt,a4paper]{amsart}
\usepackage[margin=19mm]{geometry}
\usepackage{amsmath,amssymb,mathtools}
\usepackage[scr=rsfs,cal=euler]{mathalfa}
\usepackage{xfrac}
\usepackage[unicode,hidelinks,bookmarks=false]{hyperref}
\newcommand{\ie}{i.e.\ }
\newcommand{\cf}{cf.\ }
\newcommand{\resp}{respectively\ }
\newcommand{\Subsection}{\subsection}
\providecommand{\nobreakdash}{\nobreak\hbox{-}}
\setlength{\emergencystretch}{2em}
\multlinegap0pt
\usepackage{mathtools}
\usepackage{amssymb}
\usepackage{xcolor}
\newtheorem{lemma}{Lemma}
\newtheorem{theorem}{Theorem}
\newcommand{\End}{\operatorname{End}}
\newcommand{\Hom}{\operatorname{Hom}}
\newcommand{\id}{\mathrm{id}}
\title{Logarithmic spectral correspondence for $V$--twisted Higgs bundles on punctured curves}
\author{Pradip Kumar}
\date{Source: arXiv:2603.15496v1 (CC BY 4.0)}
\begin{document}
\maketitle

\noindent\emph{Source and licence.} Logarithmic spectral correspondence for $V$--twisted Higgs bundles on punctured curves. Version arXiv:2603.15496v1 (CC BY 4.0), distributed by arXiv under the Creative Commons Attribution 4.0 International licence, http://creativecommons.org/licenses/by/4.0/, as recorded in the arXiv licence metadata for this exact version and shown on the abstract page https://arxiv.org/abs/2603.15496v1. Full licence notice: \texttt{licenses/arxiv-2603.15496-CC-BY-4.0.txt}.\par\smallskip

\noindent\textit{Editorial scope.} Counted unit: the article's theorem thm:A (source0.tex lines 555--572). The article's complete original proof of it is delivered with it (source0.tex lines 573--641, one \texttt{proof} environment of 69 lines). No mathematics is written, repaired or re-derived: the statement and the proof are the article's own lines, in the article's own notation, and no retained line is substituted or re-flowed.  Retained uncounted as the source-local context the proof needs: lines 286--288; lines 408--412; lines 422--424; lines 425--427; lines 432--443; lines 434--438; lines 534--536.  Nothing was omitted from any retained span, so the package carries no omission record.  No retained line contains a citation, so the wrapper carries no bibliography and no bibliography entry.  No retained line contains a figure environment, an image-inclusion command or drawing code, so no image is delivered and the package declares no asset dependency.  Editorial formatting only, in addition: an AMS \texttt{amsart} preamble that restates the article's own declarations and macros used by the retained lines, namely the environments \texttt{lemma}, \texttt{theorem} on separate counters, together with the standard packages those lines need.  The retained environments print in this document as Lemma 1, Theorem 1 in order of appearance, whereas the article prints the same items with the numbering of its own full document.  The complete native source of the article as distributed in the arXiv e-print for this exact version is bundled unchanged as \texttt{sources/196483/source0.tex}.
\begin{equation}\label{eq:log-hecke}
0\longrightarrow V \xrightarrow{\ \Phi\ } W(P)\xrightarrow{\ \xi\ } T_{\mathrm{tot}}\longrightarrow 0,
\end{equation}

\begin{equation}\label{eq:integrability}
\theta\wedge\theta\,=\,0
\qquad\text{in}\qquad
H^0\bigl(X,\,\, \End(E)\otimes \wedge^2V\bigr).
\end{equation}

\begin{equation}\label{eq:ThetaTheta}
\Theta\,:=\,(\id_{\End(E)}\otimes \widehat{\rho}_1)(\theta)\in H^0\bigl(X,\,\, \End(E)\otimes S(P)\bigr),
\end{equation}

\begin{equation}\label{eq:ThetaPrime}
\Theta'\,:=\,(\id_{\End(E)}\otimes \widehat{\rho}_2)(\theta)\in H^0\bigl(X,\,\, \End(E)\otimes L(P)\bigr).
\end{equation}

\begin{lemma}\label{lem:comm}
Let \(\theta\in H^0\bigl(X,\,\, \End(E)\otimes V\bigr)\) satisfy \eqref{eq:integrability}, and let \(\Theta,\,\Theta'\) be defined by \eqref{eq:ThetaTheta} and \eqref{eq:ThetaPrime}.  Then
\begin{equation}\label{eq:comm}
(\Theta'\otimes \id_{S(P)})\circ \Theta
\,=\,
(\Theta\otimes \id_{L(P)})\circ \Theta'
\end{equation}
as meromorphic homomorphisms
\[
E\longrightarrow E\otimes S(P)\otimes L(P).
\]
\end{lemma}

\begin{equation}
(\Theta'\otimes \id_{S(P)})\circ \Theta
\,=\,
(\Theta\otimes \id_{L(P)})\circ \Theta'
\end{equation}

\begin{equation}\label{eq:hecke-constraint}
\mathsf{C}_a(\Theta,\,\Theta')\,=\,0.
\end{equation}

\begin{theorem}\label{thm:A}
Let \(E\) be a holomorphic vector bundle on \(X\), and let
\[
\Theta\in H^0\bigl(X,\,\, \End(E)\otimes S(P)\bigr),\qquad
\Theta'\in H^0\bigl(X,\,\, \End(E)\otimes L(P)\bigr)
\]
be a pair of logarithmic line--twisted fields.

\begin{enumerate}
\item[(a)] \textup{(lifting)} There exists a unique section
\[
\theta\in H^0\bigl(X,\,\, \End(E)\otimes V\bigr)
\]
inducing \((\Theta,\,\Theta')\) via \eqref{eq:ThetaTheta} and \eqref{eq:ThetaPrime} if and only if the local logarithmic Hecke constraints \eqref{eq:hecke-constraint} hold for every \(a\in D_{\mathrm{tot}}\).

\item[(b)] \textup{(integrability)} Under the equivalent conditions in \textup{(a)}, the lift \(\theta\) is a \(V\)-twisted Higgs field, that is, \(\theta\wedge\theta\,=\,0\), if and only if the commutativity relation \eqref{eq:comm} holds.
\end{enumerate}
\end{theorem}

\begin{proof}
For part \textup{(a)}, tensor the exact sequence \eqref{eq:log-hecke} with \(\End(E)\).  Since \(\End(E)\) is locally free, we obtain a short exact sequence of coherent sheaves
\[
0\longrightarrow \End(E)\otimes V
\longrightarrow
\End(E)\otimes W(P)
\xrightarrow{\ \Psi\ }
\End(E)\otimes T_{\mathrm{tot}}
\longrightarrow 0,
\]
where \(\Psi\,:=\,\id_{\End(E)}\otimes \xi\).  Under the decomposition
\[
\End(E)\otimes W(P)
\,\cong\,
\bigl(\End(E)\otimes S(P)\bigr)\oplus \bigl(\End(E)\otimes L(P)\bigr),
\]
set
\[
s\,:=\,\Theta\oplus \Theta'
\in H^0\bigl(X,\,\, \End(E)\otimes W(P)\bigr).
\]
By left exactness of global sections, a lift
\[
\theta\in H^0\bigl(X,\,\, \End(E)\otimes V\bigr)
\]
with image \(s\) exists if and only if \(s\in \ker H^0(\Psi)\), that is, if and only if
\[
\Psi(s)\,=\,0\in H^0\bigl(X,\,\, \End(E)\otimes T_{\mathrm{tot}}\bigr).
\]
Such a lift is unique because the map on global sections induced by \(\Phi\) is injective.

Since \(T_{\mathrm{tot}}\) is supported on the finite set \(D_{\mathrm{tot}}\), one has a canonical decomposition
\[
H^0\bigl(X,\,\, \End(E)\otimes T_{\mathrm{tot}}\bigr)
\,=\,
\bigoplus_{a\in D_{\mathrm{tot}}}
\Hom\bigl(E_a,E_a\otimes (T_{\mathrm{tot}})_a\bigr).
\]
Under this identification, the \(a\)-component of \(\Psi(s)\) is exactly
\[
\mathsf{C}_{1,a}(\Theta)+\mathsf{C}_{2,a}(\Theta')\,=\,\mathsf{C}_a(\Theta,\,\Theta').
\]
Therefore \(\Psi(s)\,=\,0\) if and only if \eqref{eq:hecke-constraint} holds for every \(a\in D_{\mathrm{tot}}\).  This proves \textup{(a)}.

For part \textup{(b)}, assume that the equivalent conditions in \textup{(a)} hold, and let \(\theta\) be the unique lift.  Set
\[
X^\circ\,:=\,X\setminus D_{\mathrm{tot}}.
\]
Since \(T_{\mathrm{tot}}|_{X^\circ}\,=\,0\), the sequence \eqref{eq:log-hecke} restricts over \(X^\circ\) to an isomorphism
\[
\Phi|_{X^\circ}:V|_{X^\circ}\xrightarrow{\sim} W(P)|_{X^\circ}\,\cong\, W|_{X^\circ}
\,\cong\, S|_{X^\circ}\oplus L|_{X^\circ}.
\]
Under this identification, \(\theta|_{X^\circ}\) corresponds to
\[
\Theta|_{X^\circ}\oplus \Theta'|_{X^\circ}.
\]
By the local computation used in the proof of Lemma~\ref{lem:comm}, the vanishing of \(\theta\wedge\theta\) on \(X^\circ\) is equivalent to the commutativity relation \eqref{eq:comm} on \(X^\circ\).

Now \(\theta\wedge\theta\) is a holomorphic section of
\[
\End(E)\otimes \wedge^2V,
\]
and the difference of the two sides of \eqref{eq:comm} is a holomorphic section of
\[
\Hom\bigl(E,E\otimes S(P)\otimes L(P)\bigr).
\]
A holomorphic section on the smooth irreducible curve \(X\) vanishes identically if and only if it vanishes on the dense open subset \(X^\circ\).  Hence \(\theta\wedge\theta\,=\,0\) on \(X\) if and only if \eqref{eq:comm} holds on \(X\).  This proves \textup{(b)}.
\end{proof}

\end{document}
